\label{ch:design}

\subsection{Design requirements}

The production design follows six requirements derived from the vLLM runtime.

\begin{enumerate}
  \item \textbf{Semantic preservation.} The same new placement must be reached
  as in the unbatched path. Scheduling may change transfer order, but not the
  selected logical experts, destination slots, or replica assignment.

  \item \textbf{No endpoint contention inside a batch.} Each rank participates
  in at most one coalesced directed flow, whether as sender or receiver.
  Independent rank pairs should remain concurrent.

  \item \textbf{Determinism.} Every rank builds the same batch structure from
  replicated placement arrays without exchanging a separate schedule.

  \item \textbf{Low control-plane cost.} The algorithm runs repeatedly during
  serving. It must avoid an optimizer and minimize Python object construction
  and repeated function-entry overhead.

  \item \textbf{Communicator compatibility.} Batch boundaries must be visible
  to the communication backend, and NIXL's transfer context lifecycle must be
  respected.

  \item \textbf{Safe rollout.} The new behavior is disabled by default and is
  reachable only in asynchronous EPLB. Existing synchronous users retain the
  established migration path.
\end{enumerate}

\subsection{From placements to remote flows}

For each rank, the scheduler converts its old and new local slot arrays into
sets of logical expert identifiers. The difference
$H^{\mathrm{new}}_r\setminus H^{\mathrm{old}}_r$ identifies experts that rank
$r$ must receive remotely. An expert already present locally does not generate
a remote instruction even when its physical row changes.

The algorithm then indexes old source ranks only for expert identifiers that
are needed somewhere. This filtering avoids building a global inverse index
for experts irrelevant to the current migration. If several source replicas
exist, the ordered destinations are divided evenly among them using quotient
and remainder assignment. This reproduces the legacy source-balancing logic.

Finally, expert identifiers are grouped in a dictionary keyed by
$(s,d)$. One \code{MigrationFlow} contains the source rank, destination rank,
and immutable tuple of experts for that ordered pair. Direction is retained for
send/receive execution, although the contention test treats both endpoints
symmetrically.

\subsection{Greedy rank-disjoint grouping}

The communication graph has rank vertices and one coalesced edge for each
directed pair. A batch is a matching. The scheduler processes flows in
deterministic insertion order and places each in the first existing batch whose
endpoint set is disjoint. If no batch is compatible, it creates a new one.

Instead of maintaining a Python set for every batch, the implementation uses an
integer bit mask. For a flow $(s,d)$,

\begin{equation}
  m_{s,d}=(1\ll s)\;\lor\;(1\ll d).
\end{equation}

If $u_b$ is the mask of ranks already used in batch $b$, the flow fits exactly
when $m_{s,d}\land u_b=0$. On insertion, $u_b$ becomes
$u_b\lor m_{s,d}$. This is constant-time endpoint testing for practical EP
group sizes.

\begin{algorithm}[t]
  \caption{Deterministic first-fit migration batching}
  \label{alg:greedy}
\begin{lstlisting}[style=code,mathescape=true]
Input: coalesced flows P in deterministic order
batches = []; used = []
for flow p = (src, dst, expert_ids) in P:
    mask = (1 << src) | (1 << dst)
    for batch_index in range(len(batches)):
        if mask & used[batch_index] == 0:
            batches[batch_index].append(p)
            used[batch_index] |= mask
            break
    else:
        batches.append([p])
        used.append(mask)
return batches
\end{lstlisting}
\end{algorithm}

The greedy algorithm does not promise the minimum number of colors for every
graph. That is acceptable because batch count is secondary to a simple and
conservative concurrency cap. On four ranks, a batch contains at most two
independent directed flows. The rank-disjoint rule intentionally leaves
headroom rather than trying to saturate the NIC with migrations. Reciprocal
flows such as $r_0\rightarrow r_1$ and $r_1\rightarrow r_0$ are distinct
directed flows and therefore occupy separate batches, even though they involve
the same rank pair.

\subsection{Correctness invariants}

The scheduler and execution path maintain three key invariants.

\paragraph{Completeness.}
Each remote gain appears in exactly one coalesced pair and each pair is inserted
in exactly one batch. Hence every derived migration is scheduled once.

\paragraph{Conflict freedom.}
Suppose two flows in the same batch share rank $r$. Both masks contain bit $r$.
When the second flow is considered, its mask intersects the batch mask, so the
algorithm rejects that batch. This contradicts co-residence; therefore flows in
a batch have disjoint endpoints.

\paragraph{Local-copy preservation.}
Experts already available on a destination rank are excluded before remote
flow construction. The existing local copy and receive-primary metadata still
control physical-row updates. Batching modifies only the remote communication
phase.

Unit tests cover deterministic grouping, rank-pair coalescing, random
instruction completeness, exclusion of local copies, replica balancing,
multi-layer equivalence, and execution with a mock communicator. The upstream
change added the test file to vLLM's expert-parallel CI area.

\subsection{Complexity}

Let $S$ be the number of physical slots across all ranks in one layer, $M$ the
number of remote expert migrations, $F$ the number of distinct directed rank
pairs after coalescing, and $K$ the number of produced batches. Placement
scanning is $O(S)$. Building and grouping the migrations is $O(M)$ expected
time with hash tables. First-fit placement is $O(FK)$ in the worst case and
$O(F)$ storage, in addition to the placement indices.

For a fixed small number of ranks, $K$ is bounded and the practical cost is
close to linear in placement size plus migration count. More importantly, the
implementation invokes the outer scheduler once for all layers in a rebalance
cycle. If the model contains $L$ MoE layers, it still constructs $L$
independent schedules, but avoids $L$ separate profiled entry points and
repeated surrounding Python setup.

\subsection{Cycle-wide precomputation, layer-wise transfer}

An earlier design called the scheduler immediately before each layer transfer.
The final design receives two-dimensional old and new placement arrays and
computes a list of schedules---one per layer---inside one cycle-wide call. This
does \emph{not} combine expert payloads from different layers. The asynchronous
worker later iterates over layers and passes only the corresponding schedule to
the transfer function.

The distinction is important. Weight transfers and mapping activation preserve
the existing layer lifecycle, limiting memory and synchronization changes.
Only CPU schedule construction is aggregated. Qwen3-30B-A3B has 48 layers, so
one profiling range represents schedule construction for all 48 layers. The
worker waits for the new placement maps before this step; adding a few
milliseconds of CPU work once per rebalance cycle is negligible relative to
the measured cycle interval, but this claim is verified rather than assumed in
\cref{sec:profile-results}.

\subsection{Batched execution}

For a layer, the executor first creates maps from logical expert identifier to
the old and new local physical row. It then processes the ordered list of
batches. For every batch it:

\begin{enumerate}
  \item restores the communicator's transfer context for the current old
  placement and layer;
  \item posts sends for flows whose source is the local EP rank;
  \item posts receives into expert buffers for flows whose destination is the
  local EP rank; and
  \item calls \code{communicator.execute()} to preserve the batch boundary.
\end{enumerate}

For NIXL, restoring context per batch is required because execution drains the
queued operations and clears the transfer context. Calling \code{execute()}
once after all batches would submit the operations together, erasing the very
boundaries that enforce rank-disjoint concurrency. Conversely, an uninvolved
rank can have no send or receive in a given batch; the backend and existing
layer coordination govern progress.

\subsection{Configuration and failure containment}

The change adds \code{enable\_migration\_batching} to \code{EPLBConfig}, with
default \code{false}. The effective property is

\begin{equation}
  \texttt{migration\_batching\_enabled}
  =\texttt{use\_async}\land\texttt{enable\_migration\_batching}.
\end{equation}

Thus, setting the new flag cannot alter synchronous EPLB. This is a deliberate
engineering boundary. In synchronous mode, serving is stopped and serial
batches can only delay its resumption. In asynchronous mode, the operator can
enable batching when network interference is a concern. Keeping the feature
opt-in also acknowledges that the current evidence comes from one cluster and
that the optimal concurrency limit may vary with topology.

When disabled, the original path constructs local metadata, posts all remote
sends and receives, and invokes the communicator once. The feature does not
change model routing, balancing policy, or numerical output. This small impact
surface was important for upstream review and made rollback a configuration
change rather than a code revert.

\subsection{Upstream integration}

The final patch added a scheduler module, unit tests, one configuration field,
async-worker wiring, and a batched branch in the migration executor. It did not
ship the offline ILP, datasets, experimental scripts, or thesis artifacts as
runtime dependencies. The contribution was merged into vLLM on 3 October 2026
as pull request \prnumber\ \cite{vllmpr52641}. The separation between research
artifacts and production files keeps the upstream change reviewable and avoids
coupling online serving to the experimental environment.
